% Generated by Listas/make_statements.py from pratica-lecture-06.md, item 22. Do not edit.
\textbf{Source.} Muñoz-Garcia \& Toro-Gonzalez (2019), \emph{Strategy and Game Theory: Practice Exercises with Answers}, Ch. 6, Exercise 10 ``Collusion and Imperfect Monitoring'', p.~310 {[}PDF 321{]}. Reproduced verbatim (statement only; the book's worked answer follows it).



\begin{quote}

Two firms produce a differentiated product, and every firm \(i = \{1, 2\}\) faces the following linear demand:



\[p_i = a - bq_i - dq_j, \quad \text{where } a, b, d > 0 \text{ and } 0 < d < b\]



Intuitively, when \(d \to 0\) firm \(j\)'s output decisions do not affect firm \(i\)'s sales (i.e., products are heterogeneous), while when \(d \to b\), firm \(j\)'s output decisions have a significant impact on firm \(i\)'s sales, i.e., products are relatively homogeneous. The firms expect to compete against each other for infinite periods, that is, there is no entry or exit in the industry. {[}\emph{Hint}: Assume that when a firm detects deviations from collusive behavior, it retaliates by reverting to the non-cooperative Cournot (Bertrand) production level (pricing, respectively), in all subsequent periods. Prior to detection, each firm assumes that its rival follows the collusive outcome.{]}



\textbf{Part (a)} Assume that firms compete as Cournot oligopolists. What is the minimum discount factor necessary in order to sustain the collusive outcome assuming that (i) Deviations are detected immediately, and (ii) Deviations are detected with a lag of one period.



\textbf{Part (b)} Repeat part (a) of the exercise under the assumption that firms compete as Bertrand oligopolists. (You can assume immediate detection)

\end{quote}

